\chapter{The frame}
\label{ch:frameformat}

{\large\itshape Field by field, from SOF to EOF.\par}

\begin{chapterbox}{In this chapter}
\begin{itemize}
  \item The data frame's seven fields, and what each of them is for.
  \item The identifier, which is two things at once: address and priority.
  \item Why CAN has no destination address.
  \item DLC, and why it is four bits when the payload is at most eight bytes.
  \item The acknowledgement that every receiver writes into the same bit.
\end{itemize}
\end{chapterbox}

\section{The shape of the frame}

A CAN data frame with an 11-bit identifier looks like this:

\bookfigure{frame_layout}{The data frame's fields. The widths are in bits, before bit
stuffing.}{fig:frame}

Seven fields, and one boundary that matters more than the others: everything up to and including
the RTR bit is the \textbf{arbitration field}, and that is where several nodes sending at the same
time settle who gets to continue (\kapref{ch:arbitrering}).

\begin{booktable}{@{}llL@{}}
\thead{Field} & \thead{Bits} & \thead{Content} \\ \midrule
SOF & 1 & Start of frame. A dominant bit, saying that a frame begins. \\
Identifier & 11 & What the frame contains, and with it its priority. \\
RTR & 1 & Remote transmission request: dominant for a data frame. \\
IDE & 1 & Identifier extension: dominant for 11-bit identifiers. \\
r0 & 1 & Reserved, dominant. \\
DLC & 4 & Data length code: the number of data bytes, 0--8. \\
Data & 0--64 & The payload, zero to eight bytes. \\
CRC & 15 & Checksum over everything before it. \\
CRC delimiter & 1 & Recessive. \\
ACK slot & 1 & The transmitter lets go; a receiver pulls it down. \\
ACK delimiter & 1 & Recessive. \\
EOF & 7 & End of frame, seven recessive bits. \\
\end{booktable}

After EOF there is also the \textbf{intermission}, three recessive bits during which nobody may
start sending. Together with EOF that makes ten recessive bits in a row between two frames --- a
pattern that cannot occur anywhere else, because bit stuffing (\kapref{ch:stuffing}) means that five
equal bits in a row is the most that ever appears inside a frame.

\begin{aside}
\note[The arithmetic] $1 + 12 + 6 + 64 + 16 + 2 + 7 = 108$ bits for the longest data frame, before
stuffing. Bit stuffing adds at most twenty-odd, so a maximal frame is around 130 bits. At
1~Mbit/s it therefore takes about 130~µs. That figure is worth remembering: it is the answer to
``how long can a transmission go on?''.
\end{aside}

\section{The identifier}

The identifier is CAN's most distinctive design choice, and it is two things at once.

\subsection{It says what the frame contains}

It is \textbf{not} a receiver address. It describes the \emph{content}:

\begin{narrowtable}{@{}ll@{}}
\thead{Identifier} & \thead{Meaning} \\ \midrule
\code{0x301} & Engine speed \\
\code{0x302} & Coolant temperature \\
\code{0x400} & Status request \\
\end{narrowtable}

A node sending \code{0x301} does not say ``to the instrument cluster'', it says ``the engine speed
is this''. Who listens is not the sender's problem, and that has consequences that are easy to miss
at first:

\begin{itemize}
  \item \textbf{Several receivers can act on the same frame.} Both the instrument cluster and the
    gearbox controller may need the engine speed, and neither needs to know about the other.
  \item \textbf{A new node can be added without changing any old one.} It simply starts listening
    for the identifiers it cares about.
  \item \textbf{Filtering is done by the receiver.} Every node receives every frame and decides for
    itself whether it concerns it --- in software, or in the controller's filter hardware.
\end{itemize}

\subsection{It sets the priority}

At the same time the identifier is the frame's priority, and there \textbf{the lower numerical
value wins}. That follows directly from \kapref{ch:buss}: the identifier is sent MSB first, a zero
is dominant, and the dominant bit wins.

This means the address space has to be planned. An identifier is not just a name; it is a place in
a priority order that applies to the whole bus. Brakes get low numbers, cabin heating high ones.

\begin{aside}
\note[A common pitfall] Two nodes must never send the \emph{same} identifier. Two frames with the
same identifier are identical bit for bit through the whole arbitration field, so neither node
detects the other --- and as soon as their data fields differ, a bit error occurs in the middle of
the frame. Identifiers are assigned per sending node, not per message type.
\end{aside}

\subsection{Extended identifiers}

CAN 2.0B allows 29-bit identifiers. The difference shows in the IDE bit: recessive instead of
dominant, followed by another 18 identifier bits. This book sticks to 11 bits, but it is worth
knowing why the choice shows in the frame: a node that understands only 11 bits must be able to
see, early and unambiguously, that it does not understand this frame.

Note the priority consequence too: an 11-bit frame with the same leading bits always wins over a
29-bit one, because its IDE bit is dominant.

\section{RTR and remote frames}

The RTR bit distinguishes two kinds of frame:

\begin{narrowtable}{@{}lll@{}}
\thead{RTR} & \thead{Kind of frame} & \thead{Meaning} \\ \midrule
\dom{} (0) & Data frame & ``Here is the data.'' \\
\rec{} (1) & Remote frame & ``Send me the data.'' \\
\end{narrowtable}

A \textbf{remote frame} has the same identifier as the data frame it asks for, a DLC but \emph{no}
data field. The idea is that a node can ask for a value instead of waiting for it to be sent.

In practice remote frames are rarely used in modern designs. They complicate more than they add:
two nodes can answer the same request, and a node that does not understand the mechanism sees a
frame whose DLC does not match the length of its data field.

That RTR sits \emph{inside} the arbitration field, on the other hand, is not arbitrary. It means a
data frame always wins over a remote frame with the same identifier --- whoever has the data takes
precedence over whoever asks for it, which is exactly the right priority.

\section{The control field and DLC}

Three bits of administration, then four bits of length:

\begin{itemize}
  \item \textbf{IDE} --- dominant for 11-bit identifiers.
  \item \textbf{r0} --- reserved, sent dominant. It is there for future use, and that future never
    came for classical CAN.
  \item \textbf{DLC} --- four bits, value 0--8.
\end{itemize}

Four bits hold 0--15, but only 0--8 are meaningful. What happens to 9--15? In the standard they are
reserved and are to be treated as 8. In practice it is one of the points where controllers differ
--- and that is why a driver that \emph{sends} must never produce a DLC above 8, whatever the
hardware happens to do with it.

\section{The CRC field}

A fifteen-bit checksum, computed over everything from SOF up to and including the data field,
followed by a recessive delimiter. How it is computed is described in \kapref{ch:stuffing}.

The delimiter is there for a practical reason: it gives the receivers a known, recessive bit to
brace against before the ACK slot, which is the only bit in the whole frame where someone other
than the transmitter drives the bus.

\section{The ACK field}

Here something happens that has no counterpart in a home-grown protocol.

The transmitter sends the ACK slot \textbf{recessive} --- that is, it lets go. Every receiver that
has received the frame and found the CRC to be correct pulls it down \textbf{dominant} at the same
moment. The transmitter reads the bus back and sees:

\begin{narrowtable}{@{}ll@{}}
\thead{The transmitter reads} & \thead{Conclusion} \\ \midrule
\dom{} & At least one node received the frame without error \\
\rec{} & Nobody did \\
\end{narrowtable}

Three things are worth dwelling on:

\textbf{The acknowledgement costs nothing.} It is a single bit, inside the frame, and it needs no
reply frame. Compare that with an ACK message that has to be built, sent, arbitrated again and
acknowledged in its turn.

\textbf{It does not say \emph{who}.} Every receiver pulls down the same bit at the same moment. The
transmitter knows that \emph{somebody} heard, never how many or which. If that is needed, the
application has to reply with a frame of its own.

\textbf{It does not say that anyone \emph{cared}.} A node that filters the frame out as
uninteresting acknowledges it anyway, because it received it without error. The acknowledgement is
a receipt for the \emph{transfer}, not for the \emph{processing}.

\begin{simplification}
\note[Simplification] The simplified controller in \kapref{ch:kursen} \emph{lets go} of the ACK slot
as a transmitter and \emph{pulls it down} as a receiver, exactly as the standard says --- but it
never reads back what happened there. A frame that nobody acknowledges therefore completes as if
everything went well.
\end{simplification}

\section{EOF and intermission}

Seven recessive bits end the frame, and three more of intermission keep the bus idle before anyone
may start again.

Seven recessive in a row is more than bit stuffing ever allows inside a frame, which makes EOF an
unambiguous end marker: a receiver that sees seven recessive bits knows the frame is over, no
matter where it lost the thread.

This is also where error signalling happens: a node that detects an error aborts by sending six
dominant bits in a row --- a pattern that is impossible in a correct frame, and that therefore
cannot be mistaken for data. \Kapref{ch:fel} goes through it.

\section*{Exercises}

\exercise{Read the frame}{Decoding}
\label{ex:2:1}
A frame has the identifier \code{0x2A7}, DLC 5 and the data bytes \code{01 02 03 04 05}.
\begin{enumerate}[a)]
  \item How many bits is the frame, before bit stuffing?
  \item How long does it take at 1~Mbit/s, before stuffing?
  \item Which identifier would have won over it in an arbitration: \code{0x2A6} or
    \code{0x2A8}?
\end{enumerate}

\exercise{No address}{Reasoning}
\label{ex:2:2}
A system has a temperature sensor and two nodes that both need the temperature.
\begin{enumerate}[a)]
  \item How many frames does the sensor have to send?
  \item What would the same thing have looked like in a protocol with a destination address?
  \item A third node is added that also wants the temperature. What has to change in the sensor?
\end{enumerate}

\exercise{The acknowledgement}{Reasoning}
\label{ex:2:3}
A node sends a frame and reads back a dominant ACK slot.
\begin{enumerate}[a)]
  \item What does it know for certain?
  \item What does it not know?
  \item The node is alone on the bus and sends. What does it read in the ACK slot, and what should a
    standard-compliant controller do then?
\end{enumerate}

\exercise{DLC}{Reasoning}
\label{ex:2:4}
DLC is four bits, but only the values 0--8 are meaningful.
\begin{enumerate}[a)]
  \item Why was the field made four bits wide?
  \item A driver writes DLC 9 to its controller. Whose responsibility is it that this does not
    happen, and why that side in particular?
\end{enumerate}
